\documentclass[onecolumn,authoryear]{els-mrw} 

\usepackage{amsmath,amssymb,amsfonts,amsthm,makeidx,graphicx}
\usepackage{txfonts}
\usepackage{helvet}

%%Please add any additional required packages before this commented line.

\begin{document}

\chapter{Article title}\label{chap1}

\author[1]{First Author}%
\author[2]{Second Author}%

\author[1,2]{Third Author}%

\address[1]{\orgname{Name of Institute}, \orgdiv{Division or Department}, \orgaddress{Address of Institute}}
\address[2]{\orgname{Name of Institute}, \orgdiv{Division or Department}, \orgaddress{Address of Institute}}

\articletag{Chapter Article tagline: update of previous edition,, reprint..}

\maketitle


\begin{glossary}[Glossary]
\term{Europe} the model is a coherent view of capital markets data that allows users to interact with the content in a consistent manner.

\term{Primates} regardless of the source. Essentially, of sources. Properly deployed.

\end{glossary}

\begin{glossary}[Nomenclature]
\begin{tabular}{@{}lp{34pc}@{}}
AF &Assessment Factor\\
ECHA &European Chemical Agency\\
EPM &Equilibrium Partitioning Method Equilibrium Partitioning Method Equilibrium Partitioning Method Equilibrium\hfill\break Partitioning Method\\
ERA &Ecological Risk Assessment\\
HC &Hazardous Concentration\\
\end{tabular}
\end{glossary}

\begin{abstract}[Abstract]
Although the typical acquirer underperforms the market in the long term, not all acquisitions
destroy value. We search for factors that could help investors screen for valueenhancing
acquirers. For instance, cash-financed deals within the same industry involve
relatively large targets, enjoying better fortunes, and bidders who trade short-term market
reaction to the deal announcement are also important leading indicators.
\end{abstract}

\section{Level 1 heading in sentance case}\label{chap1:sec1}

The results of this research are largely equivocal since body image is such a personal
concept, however some research suggests that appearance-related issues may, to some
extent, be influenced by treatment, age-related ordemographic details. For example,
younger patients with head and neck cancer have reported higher levels of anxiety
prior to disfiguring surgery\footnote{Gases and aerosols that move into or out of the
Earth's land and water bodies and those that are lost to space are
presumed to have an inconsequential effect on the mass present.} and more concerns around body image and sexuality have
beenreported by younger patients with breast, colon, testicular, female reproductive or
lymphatic cancer than by older patients with the same disease.

Where is the internuclear distance between spins, are the gyromagnetic
The angular portyion of the dipolar Hamiltonian is described using the second rank
Legendre function which is a function of the angle uij subtended by the magnetic field.
The psychosocial impact of an altered appearance after cancer can be far-reaching and
varied. As to the (negative) relationship between profits and share prices with respect to M\&As, see profits and
share prices with respect \cite{bib1}.

Where is the internuclear distance between spins, are the gyromagnetic
The angular portyion of the dipolar Hamiltonian is described using the second rank
Legendre function which is a function of the angle uij subtended by the magnetic field.
The psychosocial impact of an altered appearance after cancer can be far-reaching and
varied. Not only do patients have to contend with a new body image, \cite{bib2} they must also
deal with their feelings towards the loss of their previous looks.

\subsection{Level 2 heading in sentance case}\label{chap1:subsec1}

The results of this research are largely equivocal since body image is such a personal
concept, however some research suggests that appearance-related issues may, to some
extent, be influenced by treatment, age-related ordemographic details. For example,
younger patients with head and neck cancer have reported higher levels of anxiety
prior to disfiguring surgery and more concerns around body image and sexuality have
beenreported by younger patients with breast, colon, testicular, female reproductive or
lymphatic cancer than by older patients with the same disease.

\subsubsection{Level 3 heading in sentance case}\label{chap1:subsubsec1}

The results of this research are largely equivocal since body image is such a personal
concept, however some research suggests that appearance-related issues may, to some
extent, be influenced by treatment, age-related ordemographic details. For example,
younger patients with head and neck cancer have reported higher levels of anxiety
prior to disfiguring surgery and more concerns around body image \cite{bib3} and sexuality have
beenreported by younger patients with breast, colon, testicular, female reproductive or
lymphatic cancer than by older patients with the same disease.

\paragraph{Level 4 heading in sentance case}

The results of this research are largely equivocal since body image is such a personal
concept, however some research suggests that appearance-related issues may, to some
extent, be influenced by treatment, age-related ordemographic details. For example,
younger patients with head and neck cancer have reported higher levels of anxiety
prior to disfiguring surgery and more concerns around body image and sexuality have
beenreported by younger patients with breast, colon, testicular, female reproductive or
lymphatic cancer than by older patients with the same disease. To be in equilibrium, the intensity
of radiation cannot be dependent on direction (i.e., radiation must be
\emph{isotropic}), and temperature cannot depend on the frequency and
direction of electromagnetic radiation \cite{bib4}.

\subparagraph{Level 5 heading in sentance case}%

The results of this research are largely equivocal since body image is such a personal
concept, however some research suggests that appearance-related issues may, to some
extent, be influenced by treatment, age-related ordemographic details. For example,
younger patients with head and neck cancer have reported higher levels of anxiety
prior to disfiguring surgery and more concerns around body image and sexuality have
beenreported by younger patients with breast, colon, testicular, female reproductive or
lymphatic cancer than by older patients with the same disease.

\subsubparagraph{Level 6 heading in sentance case}%

The results of this research are largely equivocal since body image is such a personal
concept, however some research suggests that appearance-related issues may, to some
extent, be influenced by treatment, age-related ordemographic details. 

\subsubsubparagraph{Level 7 heading in sentance case}%

The results of this research are largely equivocal since body image is such a personal
concept, however some research suggests that appearance-related issues may, to some
extent, be influenced by treatment, age-related ordemographic details.

\section{Display Math}\label{chap1:sec4}

All factors were found to be statistically significant. The goodness of fit from the
regression is also quite satisfactory. The six factors we used are capable of explaining
about 6\% of the return variability. The intercept from the regression is negative,
confirming that the average acquirer under performs the market. The long-term returns
of acquirers, however, are positively associated with the short-term market reaction,
the relative size of the target, and the valuation attractive be an indication
of a value enhancing
or a value-destroying acquisition.
\begin{align}\label{chap1:eq1}
p[m_1]+\cdots+p[m_2]=p[n_1]+\cdots+p[n_2]
\end{align}

In this chapter, we provide an overview of the theory of population games and deterministic evolutionary dynamics.  We introduce population games through a series of examples and illustrate their basic geometric properties.  We formally derive deterministic evolutionary dynamics from revision protocols, introduce the main families of dynamics---imitative/biological, best response, comparison to average payoffs, and pairwise comparison and discuss their basic properties.
\begin{align}\label{chap1:eq2}
\mathrm{{H_{2}}^{+}} + \mathrm{e}^{-} & \rightarrow  \mathrm{H} + \mathrm{H}, \\
\mathrm{HeH^{+}} + \mathrm{e}^{-} & \rightarrow  \mathrm{He} + \mathrm{H}.\label{chap1:eq3}
\end{align}

In this chapter, we provide an overview of the theory of population games and deterministic evolutionary dynamics.  We introduce population games through a series of examples and illustrate their basic geometric properties.  We formally derive deterministic evolutionary dynamics from revision protocols, introduce the main families of dynamics---imitative/biological, best response, comparison to average payoffs, and pairwise comparison and discuss their basic properties.

\begin{figure}[t]
\centering
\includegraphics[width=.4\textwidth]{blankfig}
\caption{A conservation relationship can also be written for electric charge, but in mesoscale modeling, electromagnetic effects are not considered to be dynamically or thermodynamically important on the model-resolved mesoscale. They are certainly important on  cloud and precipitation microphysics, and can therefore affect mesoscale
and larger processes, but this would need to be included through.}
\label{chap1:fig1}
\end{figure}

\section{Lists}\label{chap1:sec5}%
Deterministic evolutionary dynamics (See the \eqweblink{Equations}{chap1:eq1} and \eqref{chap1:eq2})  reflect the character of the protocols that generate them; for example, dynamics based on imitation are readily distinguished from those based on optimization \cite{bib5}.

These principles form a coupled set of relations that must be
satisfied simultaneously and that include sources and sinks in the
individual expressions.
\begin{itemize}%
\item The atmosphere is dry with no phase changes of water occurring.
\item Comparatively short time periods are involved so that radiational
heating or cooling of the air is relatively small.
\begin{itemize}%
\item conservation of motion,
\item conservation of water, and
\begin{itemize}%
\item conservation of mass,
\item conservation of heat,
\end{itemize}
\item conservation of other gaseous and aerosol materials.
\end{itemize}
\item The heating or cooling of the lowest levels of the atmosphere by
the bottom surface is of comparatively small magnitude.
\end{itemize}%%
Nevertheless, in some classes of games, dynamics derived from a variety of choice principles exhibit qualitatively similar behavior.
\begin{enumerate}%
\item The atmosphere is dry with no phase changes of water occurring.
\item Comparatively short time periods are involved so that radiational
heating or cooling of the air is relatively small.
\begin{enumerate}[a.]%
\item conservation of motion,
\item conservation of water, and
\begin{enumerate}%
\item conservation of motion,
\item conservation of other gaseous and aerosol materials.
\end{enumerate}
\item conservation of other gaseous and aerosol materials.
\end{enumerate}
\item The heating or cooling of the lowest levels of the atmosphere by
the bottom surface is of comparatively small magnitude.
\end{enumerate}%%

The expression for work in \eqweblink{Eq.}{chap1:eq3} could also have included external
work performed by such processes as \nobreak chemical reactions, phase changes,
or electromagnetism; however, these effects are not included in this
derivation of work.
\begin{unenumerate}%%
\item The atmosphere is dry with no phase changes of water occurring.
\item Comparatively short time periods are involved so that radiational
heating or cooling of the air is relatively small.
    \begin{unenumerate}%%
    \item conservation of motion. The heating or cooling of the lowest levels of the atmosphere\index{atmosphere} by
    the bottom surface is of comparatively small magnitude.
    \begin{unenumerate}%%
    \item conservation of water, and the heating or cooling of the lowest levels of the atmosphere by
    the bottom surface is of comparatively small magnitude.
    \item conservation of other gaseous and aerosol materials. The heating or cooling of the lowest levels of the atmosphere by
    the bottom surface is of comparatively small magnitude.
    \end{unenumerate}%%
\item The heating or cooling of the lowest levels of the atmosphere by
the bottom surface is of comparatively small magnitude.
    \end{unenumerate}
\item The heating or cooling of the lowest levels of the atmosphere by
the bottom surface is of comparatively small magnitude.
\end{unenumerate}%%
The ideal gas law, referred to previously, was derived from
observations of the behavior of gases at different pressures,
temperatures, and volumes \cite{bib6}.
\begin{description}
\item[A setup-time $T_c$ per pair of nuclear centers.] For interelectron repulsion
  integrals, the matrix $J_{\mu'\mu}({\mathbf{S}})$ evaluated as well.  For a
  molecule with $M$ atoms, this is done at most $M(M+1)/2$ times but
  fewer if nuclear displacements are repeated, or if symmetry
  reductions can be exploited.
\item[A setup time $T_2$ for each orbital density]  For each orbital density, the expansion
  coefficients $c_{\tau_1,\tau_2}^\tau$ are initialized. For
  one-particle integrals, nothing else needs to be done for this
  step. For the interelectron repulsion integrals.
\item[Computation time per interelectron repulsion integral $T_4$] For every integral, the final dot product is performed.
\end{description}
Investigators in the 17th and 18th
centuries found that, for a given gas, pressure times volume equals a
constant at any fixed temperature (Boyle's law) and that pressure
divided by temperature equals a constant at any fixed volume
(Charles's law). These two relations can be stated more precisely as
\begin{quote}
\quotehead{Quotehead}
The foundation for any model is a set of conservation principles. For mesoscale atmospheric models, these principles are conservation of mass, conservation of heat, conservation of motion, conservation of water, the conservation of other gaseous and aerosol materials, and an equation of state.
\source{--source}
\end{quote}
The ideal gas law, referred to previously, was derived from
observations of the behavior of gases at different pressures,
temperatures, and volumes. Investigators in the 17th and 18th
centuries found that, for a given gas, pressure times volume equals a
constant at any fixed temperature (Boyle's law) and that pressure
divided by temperature equals a constant at any fixed volume
(Charles's law). These two relations can be stated more precisely as

\begin{figure}[b]
\centering
\includegraphics[width=0.65\textwidth]{blankfig}
\caption{A conservation relationship can also be written for electric charge, but in mesoscale modeling, electromagnetic effects are not considered to be dynamically or thermodynamically important on the model-resolved mesoscale.}
\label{chap1:fig2}
\end{figure}

\section{Floats}\label{chap1:sec6}
The value of the gas constant $R$ for different gases is determined
using Avogadro's hypothesis that at a given temperature and pressure
gases containing the same number of molecules occupy the same
volume. From experimental work, for example, it has been shown that at
a pressure of 1 atm ($P_0 = 1014$ mb) and a temperature of $T_0 = 273$
K, 22.4 kliter of a gas ($V_0$) will have a mass in kilograms equal to
the molecular weight of the gas $\mu$. This quantity of gas is defined
as 1 kmol. A conservation relationship can also be written for electric charge, but in mesoscale modeling, electromagnetic effects are not considered to be dynamically or thermodynamically important on the model-resolved mesoscale. They are certainly important on  cloud and precipitation microphysics, and can therefore affect mesoscale and larger processes, but this would need to be included through~parameterizations of the microphysics. 
In this chapter, we provide an overview of the theory of population games and deterministic evolutionary dynamics.  We introduce population games through a series of examples and illustrate their basic geometric properties.  We formally derive deterministic evolutionary dynamics from revision protocols, introduce the main families of dynamics---imitative/biological, best response, comparison to average payoffs, and pairwise comparison and discuss their basic properties \cite{bib7}.

In this chapter, we provide an overview of the theory of population games and deterministic evolutionary dynamics.  We introduce population games through a series of examples and illustrate their basic geometric properties.  We formally derive deterministic evolutionary dynamics from revision protocols, introduce the main families of dynamics---imitative/biological, best response, comparison to average payoffs, and pairwise comparison and discuss their basic properties.
In this chapter, we provide an overview of the theory of population games and deterministic evolutionary dynamics.  We introduce population games through a series of examples and illustrate their basic geometric properties.  We formally derive deterministic evolutionary dynamics from revision protocols, introduce the main families of dynamics---imitative/biological, best response, comparison to average payoffs, and pairwise comparison and discuss their basic properties.
\begin{table}[t]
\TBL{\caption{Molecular weight and fractional contribution by mass
of major gaseous components of the atmosphere}\label{chap1:tab1}}
{\begin{tabular*}{\textwidth}{@{\extracolsep{\fill}}@{}lll@{}}
\toprule
\multicolumn{1}{@{}l}{\TCH{Gas}} &
\multicolumn{1}{c}{\TCH{Molecular weight\footnotemark{a}}} &
\multicolumn{1}{l}{\TCH{Fractional contribution by mass}}\\
\colrule
N$_2$ & 28.016 & 0.7551\\
O$_2$ & 32.00\phantom{6} & 0.2314\\
Ar & 39.94\phantom{6} & 0.0128\\
H$_2$O & 18.02\phantom{6} & variable\\
\botrule
\end{tabular*}}{%
\begin{tablenotes}
\footnotetext[a]{Table footnote text...}
\footnotetext{\source{Table source text...}}
\end{tablenotes}
}%
\end{table}
In this chapter, we provide an overview of the theory of population games and deterministic evolutionary dynamics.  We introduce population games through a series of examples and illustrate their basic geometric properties.  We formally derive deterministic evolutionary dynamics from revision protocols, introduce the main families of dynamics---imitative/biological, best response, comparison to average payoffs, and pairwise comparison and discuss their basic properties.
In this chapter, we provide an overview of the theory of population games and deterministic evolutionary dynamics.  We introduce population games through a series of examples and illustrate their basic geometric properties.  We formally derive deterministic evolutionary dynamics from revision protocols, introduce the main families of dynamics---imitative/biological, best response, comparison to average payoffs, and pairwise comparison and discuss their basic properties.

\begin{table}%% Coding for Non-floating table
\TBL{\caption{Table caption}}{%
\begin{tabular*}{\columnwidth}{@{\extracolsep\fill}ll@{}}
\toprule
\multicolumn{1}{@{}l}{\TCH{Gas}} &
\multicolumn{1}{c}{\TCH{Molecular weight}} \\
\colrule
N$_2$ & 28.01 \\
O$_2$ & 32.00\\
Ar & 39.94\\
H$_2$O & 18.02\\
\botrule
\end{tabular*}}{}
\end{table}

In this chapter, we provide an overview of the theory of population games and deterministic evolutionary dynamics.  We introduce population games through a series of examples and illustrate their basic geometric properties.  We formally derive deterministic evolutionary dynamics from revision protocols, introduce the main families of dynamics---imitative/biological, best response, comparison to average payoffs, and pairwise comparison and discuss their basic properties.

\section{Boxed Text}\label{chap1:sec7}%

When water vapor is included, the apparent molecular weight can be
written as where $q$ is the specific humidity or ratio of the mass of water vapor
$M$, to the mass of dry air $M_{\mathrm{d}}$. Expanding this relation This form of the ideal gas law includes the contribution of water
vapor and is often written as

Stated another way, this concept requires that the
mass into and out of an infinitesimal box must be equal to the change
of mass in the box. Such a volume is sketched,
where $\rho u|_1 \, \delta y \, \delta z$ is the mass flux into the left side
and $\rho u|_2 \,
 \delta y \, \delta z$ the mass flux out of the right side. The
symbols $\delta x$, $\delta y$, and $\delta z$ represent the
perpendicular sides of the box,  $\rho$ the density, and $u$ the velocity
component normal to the $\delta z \, \delta y$ plane.
In the Earth's atmosphere, mass is assumed (See \weblink{Boxes}{chap1:box1} and Text \weblink{Box}{chap1:box1}) to have neither sinks nor the
sources.

\begin{BoxTypeA}[chap1:box1]{Box head}
\section*{Box 1 hd}
Stated another way, this concept requires that the
mass into and out of an infinitesimal box must be equal to the change
of mass in the box. Such a volume is sketched,
where $\rho u|_1 \, \delta y \, \delta z$ is the mass flux into the left side
and $\rho u|_2 \,
 \delta y \, \delta z$ the mass flux out of the right side.

\subsection*{Box 2 hd}
The
symbols $\delta x$, $\delta y$, and $\delta z$ represent the
perpendicular sides of the box,  $\rho$ the density, and $u$ the velocity
component normal to the $\delta z \, \delta y$ plane.
In the Earth's atmosphere, mass is assumed to have neither sinks nor the
sources.

Stated another way, this concept requires that the
mass into and out of an infinitesimal box must be equal to the change
of mass in the box.

Such a volume is sketched,
where $\rho u|_1 \, \delta y \, \delta z$ is the mass flux into the left side
and $\rho u|_2 \,
 \delta y \, \delta z$ the mass flux out of the right side. The
symbols $\delta x$, $\delta y$, and $\delta z$ represent the
perpendicular sides of the box,  $\rho$ the density, and $u$ the velocity
component normal to the $\delta z \, \delta y$ plane.
In the Earth's atmosphere, mass is assumed to have neither sinks nor the
sources.

Stated another way, this concept requires that the
mass into and out of an infinitesimal box must be equal to the change
of mass in the box. Such a volume is sketched,
where $\rho u|_1 \, \delta y \, \delta z$ is the mass flux into the left side
and $\rho u|_2 \,
 \delta y \, \delta z$ the mass flux out of the right side.

 The
symbols $\delta x$, $\delta y$, and $\delta z$ represent the
perpendicular sides of the box,  $\rho$ the density, and $u$ the velocity
component normal to the $\delta z \, \delta y$ plane.
In the Earth's atmosphere, mass is assumed to have neither sinks nor the
sources.
\end{BoxTypeA}

The foundation for any model is a set of conservation principles. For mesoscale atmospheric models, these principles are conservation of mass, conservation of heat, conservation of motion, conservation of water, the conservation of other gaseous and aerosol materials, and an equation of state.

The foundation for any model is a set of conservation principles. For mesoscale atmospheric models, these principles are conservation of mass, conservation of heat, conservation of motion, conservation of water, the conservation of other gaseous and aerosol materials, and an equation of state.

\section{Enunciations}
\label{chap1:App1}
where $q_1, q_2$ and $q_3$ are defined as the ratio of the mass of the
solid, \cite{bib8} and \cite{bib9} liquid, and vapor forms of water, respectively, to the mass of
air in the same volume.
\begin{theorem}
The source-sink term $S_{q_{n}}$ refers to the
processes $\sum_{i=1}^{N} m_i\,{=}\, 1$. \eqweblink{Equation}{chap1:eq1} whereby water undergoes phase changes as well as to water
generated or lost in chemical reactions.
\end{theorem}
\begin{theorem}[Proof of Theorem]
The source-sink term $S_{q_{n}}$ refers to the
processes $\sum_{i=1}^{N} m_i\,{=}\, 1$. \eqweblink{Equation}{chap1:eq2} whereby water undergoes phase changes as well as to water
generated or lost in chemical reactions.
\end{theorem}
\begin{theorem*}
The source-sink term $S_{q_{n}}$ refers to the
processes $\sum_{i=1}^{N} m_i\,{=}\, 1$. \eqweblink{Equation}{chap1:eq3} whereby water undergoes phase changes as well as to water
generated or lost in chemical reactions.
\end{theorem*}
\begin{theorem*}[Proof of Theorem]
The source-sink term $S_{q_{n}}$ refers to the
processes $\sum_{i=1}^{N} m_i\,{=}\, 1$. \eqweblink{Equation}{chap1:eq3} whereby water undergoes phase changes as well as to water
generated or lost in chemical reactions.
\end{theorem*}
\begin{proof}
The source-sink term $S_{q_{n}}$ refers to the \cite{bib10}
processes whereby water undergoes phase changes as well as to water
generated or lost in chemical reactions.
\end{proof}
For most mesoscale
applications, chemical changes in water mass can be neglected and the
terms can be expressed as contributions owing to the following
processes:
\begin{definition}
The source-sink term $S_{q_{n}}$ refers to the
processes $\sum_{i=1}^{N} m_i\,{=}\, 1$. \eqweblink{Equation}{chap1:eq2} whereby water undergoes phase changes as well as to water
generated or lost in chemical reactions.
\end{definition}
For most mesoscale
applications, chemical changes in water mass can be neglected and the
terms can be expressed as contributions owing to the following
processes:
\begin{remark}
The source-sink term $S_{q_{n}}$ refers to the
processes $\sum_{i=1}^{N} m_i\,{=}\, 1$. \eqweblink{Equation}{chap1:eq1} whereby water undergoes phase changes as well as to water
generated or lost in chemical reactions.
\end{remark}

\begin{ack}[Acknowledgments]
The result of this approach is to produce a top-down view of tIt is a
framework that standardizes the manner in which organizations can
refer to complex data content, thereby reducing the overheads supports
an incremental approach to improving business applications and
operating efficiency. It can be applied as and when needed so the manner
in which organizations can refer to complex data content, thereby
reducing the overheads supports an incremental approach to improving
business applications and that new systems take on a standard.
\end{ack}

\seealso{article title article title}

\bibliographystyle{Harvard}
\bibliography{reference}

\end{document}

